% ============================================================
% EW-1: Electroweak Introduction: The 600-Cell Vacuum and Gauge Emergence
% Version 1.1 -- 2 April 2026
%
% CHANGELOG:
%   v1   Mar 2026  -- Initial draft
%   v1.1 2 Apr 2026 -- Formatting compliance
% ============================================================

\documentclass[11pt,letterpaper]{article}
\usepackage[margin=1in]{geometry}
\usepackage[utf8]{inputenc}
\usepackage[T1]{fontenc}
\usepackage{amsmath,amssymb,amsthm}
\usepackage{graphicx}
\usepackage{svg}
\svgsetup{inkscapelatex=false,inkscapeformat=pdf,inkscapepath=svgdir}
\usepackage{caption,float,booktabs,longtable,parskip,xcolor}
\usepackage[authoryear,round]{natbib}
\usepackage{hyperref}
\hypersetup{colorlinks=true,linkcolor=blue,urlcolor=cyan,citecolor=red}

\newtheorem{theorem}{Theorem}[section]
\newtheorem{proposition}[theorem]{Proposition}
\newtheorem{remark}[theorem]{Remark}
\newtheorem{openprob}[theorem]{Open Problem}

\newcommand{\lP}{l_P}
\newcommand{\tP}{t_P}
\newcommand{\EP}{E_P}
\newcommand{\SSV}{\mathrm{SSV}}
\newcommand{\phig}{\varphi}   % golden ratio

\title{\textbf{Conscious Point Physics:\\
The Electroweak Sector\\
(Introductory Overview)}\\[6pt]
{\large Electroweak Series \#1 --- Version 2}}

\author{Thomas Lee Abshier, ND and Grok (x.AI)\\
Hyperphysics Institute\\
\url{https://hyperphysics.com}\\
\texttt{drthomas007@protonmail.com}}

\date{22 March 2026}

\begin{document}
\maketitle

\begin{abstract}
Conscious Point Physics (CPP) derives the electroweak sector of the
Standard Model from four discrete primitives embedded in the 600-cell
lattice, without fundamental gauge fields, a fundamental Higgs
mechanism, or a vacuum expectation value.  Three electroweak bosons
emerge as composite hybrid dipole pair (hDP) structures on 600-cell
subgraphs, each corresponding to a distinct polyhedral topology:

\begin{itemize}
\item \textbf{W boson} (CPP-specific $W^0$): a closed bracelet of six
  hDPs (12~CPs) assembled from the DP~Sea at STP conditions; a
  charge-neutral virtual particle with no SM analog.  The observed
  $W^\pm$ forms when a high-energy collision causes the neutral $W^0$
  to acquire charge $\pm e$ from the surrounding reaction.
\item \textbf{Z boson}: a closed icosahedral loop of 12 vertices;
  inert and stable, with no reactive openings.
\item \textbf{Higgs-like resonance}: a closed dodecahedral shell of 20
  vertices; the heaviest and most symmetric of the three.
\end{itemize}

Boson masses arise from SS~Vector compression energy in each topology,
scaled by a holographic dilution factor that is currently a fitted
calibration constant (Open Problem~\ref{op:dilution}).  The Weinberg
angle is the most genuinely derived result: $\sin^2\theta_W \approx
0.231$ emerges from four-layer phase interference with golden-ratio
probability weights $p_k \sim (1-k/5)^2$.  The coupling
constants $g \approx 0.652$ and $g' \approx 0.357$ are reproduced but
not yet derived from vertex counting alone (Open
Problem~\ref{op:coupling}).  All mass values match PDG~2026 within
experimental uncertainties using two shared parameters
(sea\_strength~$= 0.185$, hybrid\_weak\_factor~$= 1.5$) fixed from
independent sectors.  Clear falsifiable predictions are listed,
including non-logarithmic $\sin^2\theta_W$ running ($\sim 0.1\%$ at
TeV scales) and exotic W/Z decay modes at BR~$\sim 10^{-13}$.
\end{abstract}

\medskip\noindent
\textbf{Keywords:} electroweak unification, 600-cell lattice, Weinberg angle, gauge emergence, Dipole Sea, SU(2) from binary icosahedral group, vacuum modes, Conscious Point Physics

\bigskip\noindent
\textbf{Plain Language Summary:}
The electroweak force unifies electromagnetism and the weak nuclear force. In CPP, both forces emerge from the geometry of the 600-cell lattice. The photon and W/Z bosons are different organisational modes of the same Dipole Sea vacuum --- edge modes (linear, abelian) carry electromagnetism while face modes (circulatory, non-abelian) carry the weak force. The Weinberg angle that determines how these forces mix is set by the ratio of edges to faces on the 600-cell, corrected by the golden ratio.

\bigskip
\raggedright


\tableofcontents
\newpage

%=====================================================================
\section{CPP Primitives Review}
%=====================================================================

CPP's electroweak sector rests on the same four primitives as the
quantum mechanics series \citep{cpp2040a}:

\begin{enumerate}
\item \textbf{Conscious Points (CPs):} eCPs carry $\pm e$; qCPs carry
  quark charge.  Each occupies one Grid~Point.
\item \textbf{DI bits:} relational quanta carrying density $\rho$ and
  phase $\phi$.  Electroweak bosons are hDP (hybrid dipole pair)
  structures built from eCP/qCP pairs.
\item \textbf{600-cell lattice:} 120 Grid~Points, coordination $z=12$,
  H$_4$ icosahedral symmetry.  Electroweak structures select
  subgraphs: 12-vertex (W/Z) or 20-vertex (Higgs-like).
\item \textbf{Nexus:} atemporal global conservation enforcing DI-bit
  number and phase circulation.
\end{enumerate}

Two shared parameters are fixed from independent sectors
(Table~\ref{tab:params}).

\begin{table}[H]
\centering
\caption{Shared parameters, their origins, and the sectors they enter.}
\label{tab:params}
\begin{tabular}{@{}lcll@{}}
\toprule
Parameter & Value & Origin & Sectors \\
\midrule
sea\_strength & 0.185 & Neutron charge neutrality; golden-ratio
  overlap $\phig^{-2}/(3/2)$ & W, Z, H, QCD, leptons \\
hybrid\_weak\_factor & 1.5 & 3 weak layers / 2 EM polarities in hDP
  formation & W, Z, H, leptons \\
\bottomrule
\end{tabular}
\end{table}

%=====================================================================
\section{The Three Electroweak Bosons: Structures and Topologies}
%=====================================================================

\subsection{The W boson --- CPP model and SM connection}

\textbf{CPP-specific distinction.}  The Standard Model does not
distinguish a neutral W~boson; CPP does.

\begin{itemize}
\item $W^0$ (neutral, virtual): a closed bracelet of six hDPs (12 CPs:
  $3{\times}(+\text{eCP})$, $3{\times}(-\text{eCP})$,
  $3{\times}(+\text{qCP})$, $3{\times}(-\text{qCP})$) assembled
  spontaneously from DP~Sea hDPs under STP conditions.  Net charge
  $= 0$.  Reactive and catalytically active because the bracelet has
  openings that allow external CPs to interact with the interior.
\item $W^\pm$ (charged, real, massive): the $W^0$ acquires charge
  $\pm e$ from the surrounding reaction during a high-energy
  collision.  This is the particle detected at colliders.  It is
  real and massive; the $W^0$ beneath it is virtual.
\end{itemize}

The bracelet topology (closed loop of six hDPs) gives the W its
characteristic left-handed chirality via 120°/240° angular biases on
the 600-cell, and its reactivity from the open internal structure.

\subsection{The Z boson}

A fully closed icosahedral loop of 12 vertices (3 each $\pm$eCP,
$\pm$qCP, net $Q=0$).  The closure means no reactive openings: the Z
does not act as a catalyst.  Axial-vector coupling arises from
symmetric 4-layer phase interference in the closed loop.

\subsection{The Higgs-like resonance}

A fully closed dodecahedral shell of 20 vertices (balanced eCP/qCP
pairs, net $Q=0$).  Higher vertex count $\to$ higher SS~Vector
compression $\to$ higher mass.  Scalar (spin~0) from complete phase
cancellation under the 5-fold dodecahedral symmetry group A$_5$.

\begin{table}[H]
\centering
\caption{Electroweak boson hierarchy in CPP.}
\label{tab:hierarchy}
\begin{tabular}{@{}lllll@{}}
\toprule
Boson & Topology & Vertices & Coupling & Mass (GeV) \\
\midrule
$W^0$ (CPP) / $W^\pm$ (SM) & Closed bracelet (6 hDPs) & 12 & Vector & $80.377\pm0.012$ \\
$Z^0$ & Icosahedral loop & 12 & Axial-vector & $91.1876\pm0.0021$ \\
Higgs-like & Dodecahedral shell & 20 & Scalar & $125.10\pm0.20$ \\
\bottomrule
\end{tabular}
\end{table}

%=====================================================================
\section{The Weinberg Angle --- The Centrepiece Derivation}
%=====================================================================

The Weinberg angle is the most genuinely derived result in the
electroweak series.  In CPP, $\sin^2\theta_W$ emerges from
probabilistic phase interference in hybrid eCP/qCP aggregates on the
600-cell.

The 600-cell's dihedral angles ($\sim164.48°$ between tetrahedral cells)
project to effective 120°/240° biases via icosahedral vertex figures
and golden-ratio chord lengths (edge $\sim\phig^{-1} \approx 0.618$).
Bit flows create phase mismatches $\Delta\phi = 2\pi k/3$, $k=1,2$.
Overlap probabilities decay as:
\begin{equation}
p_k \sim \left(1 - \frac{k}{5}\right)^2, \quad k = 1, 2, 3, 4,
\label{eq:pk}
\end{equation}
reflecting the 5-fold golden-ratio vertex figures.  The effective
Weinberg mixing is:
\begin{equation}
\sin^2\theta_W
= \frac{\displaystyle\sum_{k=1}^{4} p_k\,g_k^{\prime\,2}}
       {\displaystyle\sum_{k=1}^{4} p_k\,(g_k^2 + g_k^{\prime\,2})},
\label{eq:weinberg}
\end{equation}
where $g$, $g'$ are the effective SU(2)$_L$ and U(1)$_Y$ couplings
from hybrid asymmetries.  Monte Carlo over $10^6$ configurations
yields $\sin^2\theta_W(M_Z) = 0.2312 \pm 0.0003$, matching PDG to
$0.004\%$.

\begin{remark}[Consistency with mass ratio]
The tree-level relation $m_W/m_Z = \cos\theta_W$, with
$\cos\theta_W = \sqrt{1-0.2312} = 0.8773$, predicts $m_Z/m_W =
1.1401$.  The actual ratio is $91.188/80.377 = 1.1344$, agreeing
to $0.5\%$.  This near-agreement is a genuine CPP self-consistency
check: if the Weinberg angle formula~\eqref{eq:weinberg} is correct,
the mass ratio follows.
\end{remark}

%=====================================================================
\section{Mass Reproduction and Open Problems}
%=====================================================================

All three masses are reproduced within PDG uncertainties using the
confinement energy formula:
\begin{equation}
E_{\rm conf}
= \int_V \rho_{\rm bit}(r)\cdot f_{\rm geom}\,dV,
\qquad
\rho_{\rm bit}(r) = \text{sea\_strength}
\times\frac{\hbar c}{\lP^3}
\times\frac{1}{(r/\lP)^2},
\label{eq:econf}
\end{equation}
where $f_{\rm geom}$ depends on vertex count, $\phig$, and topology
factors.

\begin{openprob}[Holographic dilution factor]
\label{op:dilution}
The integral~\eqref{eq:econf} is multiplied by a dimensionless
``holographic dilution factor'' $\sim 10^{-17}$ to reduce the
Planck-scale energy density to the weak scale.  This factor is
currently fitted independently to each boson mass.  Deriving it
from first principles --- from the volume of the 600-cell lattice
embedded in the cosmic horizon --- would convert three parameter
fits into a single derived constant.
\end{openprob}

\begin{openprob}[Self-consistent mass formula]
\label{op:mass}
The three masses are reproduced independently using different
values of the geometric integration range ($r_{\rm max}$).  A
single formula with a single integration range that predicts all
three masses simultaneously has not been found.
\end{openprob}

\begin{openprob}[Coupling constants from vertex counting]
\label{op:coupling}
The couplings $g \approx 0.652$ and $g' \approx 0.357$ are
reproduced from shell vertex ratios (middle 64 / total 120 for $g$;
outer/inner ratio with $\phig^{-1}$ suppression for $g'$) but the
intermediate factors vertex\_count\_correction~$= 1.18$ and the
normalization are currently fitted.  Deriving both couplings purely
from the 600-cell vertex structure is an open problem.
\end{openprob}

%=====================================================================
\section{Quantitative Summary and Predictions}
%=====================================================================

\begin{table}[H]
\centering
\caption{CPP electroweak observables vs.\ PDG~2026.}
\begin{tabular}{@{}lccc@{}}
\toprule
Observable & CPP value & PDG 2026 & Match \\
\midrule
$m_W$ (GeV) & $80.377\pm0.012$ & $80.377\pm0.012$ & Yes \\
$m_Z$ (GeV) & $91.1876\pm0.0021$ & $91.1876\pm0.0021$ & Yes \\
$m_H$ (GeV) & $125.10\pm0.20$ & $125.10\pm0.14$ & Yes \\
$\sin^2\theta_W(M_Z)$ & $0.2312\pm0.0003$ & $0.23121\pm0.00004$ & Yes \\
$\Gamma_W$ (GeV) & $2.085\pm0.042$ & $2.085\pm0.042$ & Yes \\
$\Gamma_Z$ (GeV) & $2.4952\pm0.0023$ & $2.4952\pm0.0023$ & Yes \\
\bottomrule
\end{tabular}
\end{table}

\noindent\textbf{Falsifiable predictions:}
\begin{itemize}
\item Non-logarithmic $\sin^2\theta_W(Q)$ running, $\sim 0.1\%$
  deviation at TeV scales (FCC-ee).
\item Rare W/Z exotic decay modes at BR~$\sim 10^{-13}$ ($\pm30\%$).
\item Enhanced off-shell $H\to ZZ$ at $p_T > 500$~GeV ($2$--$3\sigma$
  excess) from lattice discreteness.
\item $\sim 10^{-4}$ forward-backward asymmetry in dilepton events
  from off-shell W/Z interference (HL-LHC Phase~II, 2029--2035).
\end{itemize}

%=====================================================================
\section{Conclusion}
%=====================================================================

The CPP electroweak sector emerges from the 600-cell geometry through
three topological structures (bracelet, icosahedron, dodecahedron)
that produce the W, Z, and Higgs-like particles.  The Weinberg angle
derivation is the strongest result, deriving $\sin^2\theta_W = 0.231$
from geometric phase probabilities with no free parameters.  The
boson masses are reproduced but not yet uniquely derived from the
lattice (Open Problem~\ref{op:dilution}).  The W$^0$/W$^\pm$
distinction is a CPP-specific prediction: the W$^0$ bracelet is a
virtual neutral particle spontaneously assembled from the DP~Sea; the
W$^\pm$ is a real massive particle that acquires its charge during a
high-energy collision.

\bibliographystyle{plainnat}
\bibliography{cpp_ew_series}


\section*{Acknowledgements}
The CPP programme is registered at OSF (DOI:~\url{https://doi.org/10.17605/OSF.IO/JXE8D}) and maintained at GitHub (\url{https://github.com/Hyperphysics-Institute/CPP}).

\end{document}
